\documentclass[12pt]{article}
\usepackage[utf8]{inputenc}
\usepackage[T1]{fontenc}
\usepackage{hyperref}
\title{Gbe, Waci Language Resources for K-12 Curriculum Development}
\author{Compiled by Claude Sonnet 5 (claude-sonnet-5) for Andr\'as Kornai}
\date{}
\begin{document}
\maketitle

\section{Introduction}
Waci (Ouatchi, Watyi; ISO 639-3 wci) is a Gbe language of Togo and Benin. It belongs to a dialect continuum that includes Ewe and Gen (Mina). Speaker counts vary: an older source gives 480,000 (1991--1993) and Wikipedia gives 920,000 (2019--2021).

\section{Primary Educational Resources}
No primers or school materials were found in the search results.

\section{Linguistic and Reference Materials}
The following are cited in the Wikipedia article. The search results did not show whether they can be read online, so access should be checked.
\begin{itemize}
\item Faton (2018), \emph{Waci Speakers in Togo and Benin: A Sociolinguistic Survey}, SIL Electronic Survey Reports 2018-003.
\item Lokossa and Capo (1978), \emph{Pr\'ecis grammatical du Gengbe et du Wacigbe}.
\item Capo (1991), \emph{A Comparative Phonology of Gbe}.
\item Anson (1972), \emph{Description synchronique d'un parler Ewe: Le Watyi}.
\end{itemize}
Source for these citations: \url{https://en.wikipedia.org/wiki/Waci_language}

\section{Religious and Community Resources}
ScriptureEarth lists Bible audio for Waci from Wycliffe Benin and Wycliffe Togo: \url{https://www.scriptureearth.org/00eng.php?iso=wci}

\section{Digital and Audio Resources}
The same ScriptureEarth page lists audio recordings from Global Recordings Network: \url{https://www.scriptureearth.org/00eng.php?iso=wci}

\section{Curriculum Development Context}
Waci is a stable language with a large speaker base. No school programme or published curriculum was found. Related Gbe languages such as Ewe have been written for much longer, which may help orthography work, but this was not examined here.
\end{document}